% =============================================================================
%  arbiter_design_report.tex
%
%  Compile with:  pdflatex arbiter_design_report.tex   (run twice for the ToC)
%
%  Packages used are all standard TeX Live / MiKTeX; no external assets.
% =============================================================================
\documentclass[11pt,a4paper]{article}

\usepackage[T1]{fontenc}
\usepackage[utf8]{inputenc}
\usepackage[margin=2.4cm]{geometry}
\usepackage{booktabs}
\usepackage{longtable}
\usepackage{array}
\usepackage{amsmath}
\usepackage{xcolor}
\usepackage{listings}
\usepackage{enumitem}
\usepackage{fancyhdr}
\usepackage{caption}
\usepackage[hidelinks]{hyperref}

% ---- signal / identifier macro: avoids escaping every underscore -------------
\newcommand{\sig}[1]{\texttt{\detokenize{#1}}}

% ---- listings style ---------------------------------------------------------
\definecolor{cmtgray}{gray}{0.42}
\definecolor{kwblue}{RGB}{20,60,140}
\definecolor{bgshade}{gray}{0.97}

\lstdefinestyle{sv}{
  language=Verilog,
  basicstyle=\ttfamily\footnotesize,
  keywordstyle=\color{kwblue}\bfseries,
  commentstyle=\color{cmtgray}\itshape,
  backgroundcolor=\color{bgshade},
  frame=single,
  rulecolor=\color{gray!40},
  breaklines=true,
  showstringspaces=false,
  columns=fullflexible,
  xleftmargin=6pt,
  framexleftmargin=6pt,
  morekeywords={logic,always_ff,always_comb,localparam,generate,endgenerate,
                genvar,assign,typedef,enum,unique,posedge,negedge,parameter}
}
\lstset{style=sv}

% ---- headers ----------------------------------------------------------------
\pagestyle{fancy}
\fancyhf{}
\lhead{\small AMBA Arbiter --- Design \& Timing Closure Report}
\rhead{\small\thepage}
\renewcommand{\headrulewidth}{0.4pt}

\title{\vspace{-1.5cm}\textbf{An AMBA-Style Four-Master Bus Arbiter}\\[4pt]
       \large Design, Verification, and Timing Closure\\[2pt]
       \normalsize From $-12.426$\,ns WNS to a Board-Ready Bitstream}
\author{Astik Thukral}
\date{\today}

\begin{document}
\maketitle
\thispagestyle{fancy}

\begin{abstract}
\noindent
This report documents a synthesisable SystemVerilog bus arbiter for four AMBA-style
masters, combining round-robin selection, static priority, nested tie-breaking,
absolute lock override, and recovery from abandoned grants. It describes the
signal flow and the mechanism behind each behavioural claim, then gives a
detailed account of the timing-closure effort: the reported Worst Negative Slack
moved from $-12.426$\,ns with eleven failing endpoints to a fully constrained,
timing-met implementation at $+1.610$\,ns. The most useful finding was that
roughly 95\,\% of the original violation had nothing to do with the arbiter's
logic. The report closes with the engineering practices that made that
diagnosis possible.
\end{abstract}

\tableofcontents
\bigskip

% =============================================================================
\section{Overview and Signal Flow}
% =============================================================================

The arbiter grants a shared bus to one of four masters. Its parameters are
\sig{NUM_MASTERS}~$=4$, \sig{PRI_WIDTH}~$=3$ (priority values $0\ldots7$), and
\sig{GRACE_W}~$=4$ (the abandoned-grant window, in cycles).

\subsection{Interface}

\begin{center}
\begin{tabular}{@{}llp{8.6cm}@{}}
\toprule
\textbf{Signal} & \textbf{Dir / Width} & \textbf{Meaning} \\
\midrule
\sig{clk}          & in, 1  & System clock, 100\,MHz \\
\sig{rst_n}        & in, 1  & Active-low reset; asynchronous assertion \\
\sig{REQ}          & in, 4  & Master $i$ is requesting the bus \\
\sig{PRI}          & in, 12 & Packed priorities; master $i$ occupies \sig{PRI[i*3 +: 3]} \\
\sig{LOCK}         & in, 4  & Master $i$ requests an atomic multi-beat hold \\
\sig{VALID_XFER}   & in, 4  & Master $i$ has actually begun its transfer \\
\sig{COMPLETE}     & in, 4  & Master $i$ has finished its transfer \\
\midrule
\sig{GRANT}        & out, 4 & One-hot grant \\
\sig{slot_complete}& out, 1 & The shared completion handshake (observability) \\
\sig{slot_type}    & out, 1 & 0 = round-robin slot, 1 = priority slot \\
\bottomrule
\end{tabular}
\end{center}

\subsection{Data flow through one arbitration}

The design is entirely combinational from \sig{REQ}/\sig{PRI} to \sig{GRANT};
all state advances only on completion.

\begin{enumerate}[leftmargin=1.5em,itemsep=3pt]
\item \textbf{Round-robin candidate.} \sig{rr_pointer_unit} feeds the live
      \sig{REQ} vector and its pointer \sig{rr_ptr} into a shared
      \sig{rotate_mask_encoder}, which rotates the vector so that \sig{rr_ptr}
      maps to index~0, priority-encodes lowest-index-wins, and rotates the
      single result bit back. The output \sig{rr_grant} is the closest
      requester at or after the pointer in cyclic order.

\item \textbf{Priority candidate.} \sig{priority_comparator} evaluates
      \emph{requesting masters only} through a balanced
      $O(\log_2 N)$-deep pairwise reduction tree. Each node carries
      $(\textit{req}, \textit{pri}, \textit{mask})$ for its subrange; merging
      keeps the higher-priority side outright, or unions the two masks on an
      exact tie. The root yields \sig{tied_mask} --- the set of requesters
      holding the maximum priority --- and \sig{is_tie}.

\item \textbf{Tie resolution.} \sig{tie_break_rr_unit} runs a
      \emph{second, independent} \sig{rotate_mask_encoder} over
      \sig{tied_mask} using its own pointer \sig{tie_ptr}, producing
      \sig{tie_grant}.

\item \textbf{Slot selection.} \sig{slot_type} chooses between the two
      candidates:
      \sig{arb_grant = (slot_type == RR) ? rr_grant : tie_grant}.

\item \textbf{Lock override.} \sig{lock_ctrl} forces the grant to a latched
      holder whenever a lock is engaged:
      \sig{eff_grant = locked ? lock_grant_vec : arb_grant}.
      This has absolute precedence over both pointers and over priority.

\item \textbf{Completion.} \sig{genuine_complete} is the granted master's
      \sig{COMPLETE}, suppressed while that master still asserts \sig{LOCK}.
      \sig{grace_window_timer} independently asserts \sig{forced_complete} if a
      granted master never asserts \sig{VALID_XFER} within \sig{GRACE_W} cycles.
      \sig{completion_merge} ORs the two into the single signal
      \sig{slot_complete}.

\item \textbf{State advance.} On \sig{slot_complete} --- and only then --- the
      granting pointer advances, \sig{slot_type} toggles, any lock releases,
      and the grace counter resets.
\end{enumerate}

\begin{center}
\ttfamily\small
\begin{tabular}{@{}l@{}}
REQ ----+--> rr\_pointer\_unit ---> rr\_grant ----+ \\
~~~~~~~~|                                         | \\
~~~~~~~~+--> priority\_comparator --> tied\_mask   | \\
PRI ----+~~~~~~~~~~~~~~~~~~~~~~~~~~~~~~~|~~~~~~~~ | \\
~~~~~~~~~~~~~~~~~~~~~~~~~~~~~~~~~~~~~~~~v~~~~~~~~ | \\
~~~~~~~~~~~~~~~~~~~~~~~~~~~ tie\_break\_rr\_unit --+--> [slot\_type mux] \\
~~~~~~~~~~~~~~~~~~~~~~~~~~~~~~~~~~~~~~~~~~~~~~~~~~~~~~~~~~~~~|          \\
~~~~~~~~~~~~~~~~~~~~~~~~~~~~~~~~~~~~~~~~~~~~~~~~~~~~~~~~~~~~~v          \\
LOCK --> lock\_ctrl -------------------------> [lock override] --> GRANT \\
~~~~~~~~~~~~~~~~~~~~~~~~~~~~~~~~~~~~~~~~~~~~~~~~~~~~~~~~~~~~~|          \\
COMPLETE / VALID\_XFER --> grace\_window\_timer --+-> completion\_merge     \\
~~~~~~~~~~~~~~~~~~~~~~~~~~~~~~~~~~~~~~~~~~~~~~~~~~~~~~~v                \\
~~~~~~~~~~~~~~~~~~~~ slot\_complete --> pointers, slot\_type, lock, grace \\
\end{tabular}
\end{center}

% =============================================================================
\section{How Each Claim Is Achieved}
% =============================================================================

\subsection{Fairness: round-robin without bias}

\sig{rotate_mask_encoder} is written once and instantiated twice, so the main
pointer and the tie-break pointer are structurally identical by construction.
Rotation uses modulo arithmetic on indices rather than bit shifts, so the design
remains correct for any \sig{NUM_MASTERS}, not only powers of two.

A pointer advances \emph{past the actual winner}, not merely by incrementing
itself. With a sparse request vector the winning index can differ from the
pointer, and incrementing the pointer alone could park it back on the same
master:

\begin{lstlisting}
rr_ptr <= (winner_idx == NUM_MASTERS-1) ? '0 : winner_idx + 1'b1;
\end{lstlisting}

\subsection{Priority without starvation}

Static priority alone starves low-priority masters indefinitely.
\sig{slot_alternation_ctrl} holds a single bit that toggles
$\text{RR}\rightarrow\text{PRIORITY}\rightarrow\text{RR}$ on every
\sig{slot_complete}, independently of which master either pointer selects.

Because round-robin slots occur every second slot, and \sig{rr_ptr} cycles
through all $N$ masters, every requesting master is guaranteed a slot within
$2N$ slots --- eight, for four masters --- \emph{regardless of how badly it
loses on priority}. With priorities $\{7,1,1,1\}$ and all four requesting, the
observed pattern is

\begin{center}\ttfamily
M0,\ M3,\ M1,\ M3,\ M2,\ M3,\ M3,\ M3,\ M0,\ \ldots
\end{center}

\noindent
M3 takes five of every eight slots; M0, M1 and M2 each take one. None
disappears.

\subsection{Tie-breaking within the maximum set}

When several requesters share the maximum priority, \sig{tied_mask} is a strict
subset of the request vector, and \sig{tie_ptr} rotates \emph{within that subset
only}. Masters below the maximum reach the bus solely on round-robin slots.
The tie pointer is left untouched when a unique maximum exists, preserving the
state independence the specification relies on.

\subsection{Lock: atomic multi-beat sequences}

\sig{lock_ctrl} engages when the effective granted master asserts \sig{LOCK} and
the slot is not already completing. While held:

\begin{itemize}[leftmargin=1.5em,itemsep=2pt]
\item \sig{lock_grant_vec} forces \sig{GRANT} to the holder, so no
      re-arbitration occurs from any source, including a higher-priority
      requester arriving mid-lock;
\item \sig{freeze} holds \sig{slot_alternation_ctrl}, so the entire sequence
      consumes exactly one slot;
\item \sig{genuine_complete} is suppressed, so intermediate beats do not end the
      slot. The sequence ends only when \sig{LOCK} drops and \sig{COMPLETE}
      arrives.
\end{itemize}

Release occurs only on \sig{slot_complete}. Because the engagement path runs
only while \sig{lock_held == 0}, a master leaving lock cannot re-engage in the
same cycle: it must win the next slot fairly. This is the ``no privileged
re-entry'' rule.

\subsection{Abandoned-grant recovery}

\sig{grace_window_timer} is gated on \sig{VALID_XFER}, not on \sig{LOCK}. Once a
grant is issued, a fixed \sig{GRACE_W}-cycle window opens:

\begin{itemize}[leftmargin=1.5em,itemsep=2pt]
\item if \sig{VALID_XFER} asserts within the window, \sig{started} latches
      permanently and the check disengages --- a slow but real transfer is never
      truncated;
\item if it does not, \sig{forced_complete} pulses and the slot advances exactly
      as on a normal completion.
\end{itemize}

\subsection{One definition of completion}

\sig{completion_merge} is deliberately the only place ``a slot ended'' is
defined:

\begin{lstlisting}
assign slot_complete = genuine_complete | forced_complete;
\end{lstlisting}

Both pointers, the alternation counter, the lock controller and the grace timer
key off exactly this signal. There is therefore no way for the lock to release
on a cycle where the pointer did not advance --- a class of fairness bug that
independent per-block completion logic would permit.

% =============================================================================
\section{Known Limitations}
% =============================================================================

Stating these precisely is part of the design, not an admission against it.

\paragraph{Turn-based, not time-based.}
The guarantee bounds the \emph{number of slots} a master waits, not wall-clock
latency. If the bus owner's transfer takes $T$ cycles, everyone else waits
proportionally. Turn fairness is not bandwidth fairness: with priorities
$\{7,1,1,1\}$ and a dominant master whose transfers are $10^6$ cycles against
others' 10 cycles, M3 receives five of eight \emph{turns} but 99.9994\,\% of the
\emph{bandwidth}. Production interconnect adds bandwidth regulation for exactly
this reason. In practice AHB bounds burst length to sixteen beats, which is what
makes turn-based fairness sufficient.

\paragraph{An indefinitely held lock starves everyone.}
If a master asserts \sig{LOCK} and never releases it, \sig{genuine_complete}
stays suppressed, \sig{slot_complete} never fires, and the bus is held forever.
This is starvation in the strict sense. A lock timeout is the single
highest-value addition the design could take.

\paragraph{A started-but-never-completed transfer hangs the arbiter.}
\sig{forced_complete} is gated on \sig{!started} by design, so once
\sig{VALID_XFER} has been seen the grace window can never fire again. The
arbiter cannot distinguish ``slow'' from ``dead''. Given a choice between
occasionally hanging and occasionally truncating a live transaction, the design
chose to hang --- truncation without protocol support corrupts data. Closing
this needs a transaction watchdog \emph{and} an abort channel to the master,
which is essentially AHB \textsc{split}/\textsc{retry}. Without the notification
path, a watchdog alone trades a hang for silent corruption.

\paragraph{Undocumented assumption.}
The module header states that \sig{REQ}/\sig{PRI} remain stable within a
transfer. It does not state the load-bearing companion: \emph{a master that
starts a transfer will finish it}.

% =============================================================================
\section{Verification}
% =============================================================================

\subsection{Strategy}

Unit testbenches exist for every submodule, plus an integration bench for the
arbiter. For the board wrapper, checking is \emph{property-based} rather than
vector-based: exact grant sequences are brittle against the relative phase of
the two pointers, whereas properties hold for every legal run and catch classes
of bug that hand-written vectors miss.

\begin{center}
\begin{tabular}{@{}llp{7.4cm}@{}}
\toprule
\textbf{ID} & \textbf{Property} & \textbf{Catches} \\
\midrule
P1 & \sig{GRANT} is one-hot or zero & encoder and mux faults \\
P2 & never grant a non-requester & request-handling faults (suppressed while
     locked, where override is correct) \\
P3 & \sig{COMPLETE} is a single-cycle pulse & a held completion would
     immediately re-complete the next transfer \\
P4 & \sig{slot_type} toggles once per \sig{slot_complete}, never otherwise &
     the alternation contract, in both directions \\
P5 & \sig{GRANT} is stable within a slot & mid-transfer re-arbitration \\
-- & every requester served within $2N$ slots & \textbf{the design's headline
     fairness guarantee, checked directly} \\
\bottomrule
\end{tabular}
\end{center}

The starvation checker is the important one: it runs continuously across every
scenario rather than as a spot check, so any fairness regression anywhere in the
run is caught.

\subsection{Scenarios}

Nine scenarios: the four priority-preset demonstrations; lock; abandonment;
single-step mode; and request stability under a mid-transfer input change.
Two deserve comment.

\emph{Abandonment} carries no explicit assertion and needs none --- if the grace
window fails to fire, the wait for further slots never returns and the watchdog
trips. The hang \emph{is} the failure signal.

\emph{Button press} models mechanical bounce on both press and release. Without
it the debounce logic is never exercised and one would merely be testing a wire.

\subsection{Results}

\begin{center}
\begin{tabular}{@{}llll@{}}
\toprule
\textbf{Bench} & \textbf{Level} & \textbf{Slots} & \textbf{Result} \\
\midrule
\sig{tb_amba_arbiter_top}                & RTL, arbiter core   & --- & pass \\
\sig{tb_amba_arbiter_demo_top}           & RTL, full wrapper   & 88  & 0 errors \\
\sig{tb_amba_arbiter_demo_top_postimpl}  & implemented netlist & 88  & 0 errors \\
\bottomrule
\end{tabular}
\end{center}

The post-implementation run matched the RTL run \emph{slot for slot} across all
88 transactions, confirming that synthesis and implementation preserved
behaviour exactly.

\subsection{Two testbench techniques worth recording}

\paragraph{Parameterise for simulation.}
\sig{TICK_DIV} and \sig{DEBOUNCE_CYCLES} are \sig{parameter}, not
\sig{localparam}, purely so a testbench can shrink them. At hardware values one
transaction spans $5\times10^7$ cycles; at simulation values, twenty.

\paragraph{Netlist simulation needs a port-only bench.}
Synthesis renames, merges and absorbs internal nets, so hierarchical references
fail to bind:
\begin{lstlisting}
ERROR: [VRFC 10-2991] 'grant' is not declared under prefix 'dut'
\end{lstlisting}
Parameters also cannot be overridden in a netlist --- the values are baked into
the gates at synthesis time, so the netlist must be \emph{built} small via
synthesis generics. The port-only bench recovers slot boundaries from
\sig{slot_type} toggling on the LED outputs, since one toggle is exactly one
completed slot.

% =============================================================================
\section{The Timing-Closure Effort}
% =============================================================================

\subsection{The starting point}

The first routed report on \sig{xc7a100tcsg324-1}:

\begin{center}
\begin{tabular}{@{}lrrr@{}}
\toprule
& \textbf{WNS} & \textbf{TNS} & \textbf{Failing endpoints} \\
\midrule
Setup & $-12.426$\,ns & $-65.627$\,ns & 11 of 30 \\
Hold  & $+0.076$\,ns  & $0.000$\,ns   & 0 of 30 \\
\bottomrule
\end{tabular}
\end{center}

For an arbiter serving four masters --- 84 LUTs and 19 flip-flops --- this is an
implausible result, and that implausibility was the first useful signal.

\subsection{Diagnosis: read the paths, not the summary}

Every one of the eleven failing endpoints touched a package pin. \textbf{There
was not a single flip-flop-to-flip-flop setup violation in the design.}

The decisive path had one logic level:

\begin{center}
\begin{tabular}{@{}llrl@{}}
\toprule
\textbf{Source} & \textbf{Destination} & \textbf{Slack} & \textbf{Levels} \\
\midrule
\sig{PRI[2]} & \sig{slot_complete} & $-12.426$ & 10 (IBUF + 8 LUT + OBUF) \\
\sig{u_sa/slot_type_reg/C} & \sig{slot_type} & $-5.630$ & \textbf{1 (OBUF only)} \\
\bottomrule
\end{tabular}
\end{center}

A path containing an output buffer and \emph{no logic at all} missed by
5.63\,ns. No amount of RTL restructuring can fix that, which located the problem
outside the design.

\subsection{Root cause 1: an unsurvivable I/O budget}

The constraint file used placeholder budgets of 40\,\% of the period in each
direction:
\[
10.000 - 4.000_{\text{in}} - 4.000_{\text{out}} - 0.035_{\text{unc}}
  = 1.965\ \text{ns available inside the chip.}
\]
The fixed buffer costs on a $-1$ speed grade part are
$\text{IBUF}\approx1.00$\,ns and $\text{OBUF}\approx2.64$\,ns, totalling
3.64\,ns \emph{before a single LUT is placed}. A pad-to-pad path was over budget
by 1.7\,ns with zero logic in it.

The arbiter's own contribution to the 18.39\,ns arrival time was
$8 \times 0.124 = 0.99$\,ns --- about 5\,\%.

\subsection{Root cause 2: a pblock that never applied}

A placement region had been added on the theory that spread placement was
inflating routing. The implementation log showed it had been silently rewritten:

\begin{lstlisting}
INFO: [Vivado 12-3520] Assignment of '...' to a pblock 'pblock_arbiter'
means that all children of 'u_sa, u_gr, u_lk, u_rr, and u_tb' are in the
pblock. Changing the pblock assignment to 'u_sa, u_gr, u_lk, u_rr, u_tb'.
\end{lstlisting}

The cell list came from
\sig{[get_cells -hierarchical -filter {PRIMITIVE_LEVEL==LEAF}]}, which swept up
all 36 IBUF/OBUF cells and the BUFG. Those cannot be placed in a SLICE-only
range, so the region became unsatisfiable. The placed result landed at
\sig{SLICE_X0Y26..X3Y50} --- entirely outside the requested
\sig{SLICE_X0Y0:SLICE_X9Y19} --- with routing still at 67--80\,\% of path delay.

\textbf{Lesson: a constraint that produces no error is not necessarily a
constraint that took effect.} Only the log revealed it.

\subsection{Root cause 3: the wrong abstraction}

An arbiter is not a chip. Its \sig{REQ}/\sig{GRANT}/\sig{PRI} wires connect to
master IP inside the same fabric and never touch a package pin. Timing it
pad-to-pad measured the Artix-7 I/O ring, not the design.

A related asymmetry compounded it: the clock reached the flip-flops through
IBUF~$+$~BUFG with \textbf{4.573\,ns of insertion delay}, and on a
flip-flop-to-port path that delay is charged entirely to the launch side with
nothing on the capture side to compensate. That alone accounts for the failing
\sig{slot_type} path.

\subsection{Root cause 4: 36 unconstrained I/O}

With no \sig{PACKAGE_PIN} constraints, Vivado auto-assigned all pins around the
die perimeter, and the placer dragged 84 LUTs across 25 CLB rows to reach them.
Per-hop routing rose to 0.5--1.5\,ns against a typical 0.1--0.3\,ns.

\subsection{Two RTL findings}

\paragraph{A redundant multiplexer (14 LUTs, one logic level).}
The priority-slot grant selected between two candidates:
\begin{lstlisting}
assign prio_grant = is_tie ? tie_grant : direct_grant;
\end{lstlisting}
But \sig{rotate_mask_encoder} returns the single set bit of a \emph{one-hot}
mask for \emph{any} pointer value --- it rotates that lone bit somewhere,
priority-encodes it as the only candidate, and rotates it back where it started.
So on the unique-maximum case, where \sig{tied_mask} is already one-hot,
\sig{tie_grant} was computing exactly the vector \sig{direct_grant} did. The
entire selection was redundant:
\begin{lstlisting}
assign prio_grant = tie_grant;
\end{lstlisting}
This removed the multiplexer, \sig{unique_valid}, \sig{direct_grant_vec} and two
\sig{\$countones} reductions: $-14$ LUTs ($-17\,\%$) and one logic level.
Simulation results were bit-identical, confirming the equivalence argument
empirically.

\paragraph{A cheaper tie test.}
\sig{\$countones(tied_mask) >= 2} was replaced by
\sig{|(tied_mask \& (tied_mask - 1))}. At $N=4$ this changed neither logic depth
nor LUT count --- Vivado had already folded the population count --- but it
altered LUT packing enough to gain $0.225$\,ns. Re-running the identical design
produced a byte-identical result, proving the gain was reproducible rather than
placement noise. The rewrite is retained because it scales; the credit belongs
to packing, not to depth.

\subsection{The decisive measurement: logic versus routing}

Every critical path in this design sits at 13--15\,\% logic and 85--87\,\%
routing:

\begin{center}
\begin{tabular}{@{}lrrr@{}}
\toprule
\textbf{Path (OOC, post-optimisation)} & \textbf{Total} & \textbf{Logic} & \textbf{Route} \\
\midrule
\sig{PRI[10]} $\rightarrow$ \sig{slot_complete} & 7.544\,ns & 0.992\,ns (13.2\,\%) & 6.552\,ns (86.9\,\%) \\
\bottomrule
\end{tabular}
\end{center}

That ratio is the single most diagnostic number in the whole effort. It predicted
what happened next: removing a full LUT level bought only $0.033$\,ns, because
the router handed most of it straight back as extra wire delay. \textbf{Deleting
every LUT on the path would recover under 1\,ns.} Further logic optimisation had
nowhere to go, and the correct response was to stop.

\subsection{The progression}

\begin{center}
\begin{tabular}{@{}lrrrr@{}}
\toprule
\textbf{Stage} & \textbf{WNS} & \textbf{TNS} & \textbf{Failing} & \textbf{LUTs} \\
\midrule
Original (4\,ns I/O budgets, pblock)      & $-12.426$ & $-65.627$ & 11 / 30 & 84 \\
I/O max $4\rightarrow1$\,ns, pblock removed & $-5.192$ & $-22.256$ & 6 / 30  & 84 \\
Out-of-context (no IBUF/OBUF/BUFG)        & $+0.115$  & $0.000$   & 0 / 30  & 83 \\
\quad + cheaper \sig{is_tie}              & $+0.340$  & $0.000$   & 0 / 30  & 83 \\
\quad + redundant mux removed             & $+0.373$  & $0.000$   & 0 / 30  & \textbf{69} \\
Ported to \sig{xc7a35tcpg236-1}           & $+0.421$  & $0.000$   & 0 / 30  & 67 \\
\midrule
Board wrapper, pad level, fully constrained & $\mathbf{+1.610}$ & $\mathbf{0.000}$ & \textbf{0 / 277} & 168 \\
\bottomrule
\end{tabular}
\end{center}

\subsection{Why the final number is larger, and legitimate}

The last row is not the same measurement made kinder. The arbiter was wrapped in
a board-level design containing a reset synchroniser, an input synchroniser, a
behavioural master model, a priority-table selector, a tick generator, a
debouncer and a display stage. The arbiter's \sig{REQ}/\sig{GRANT} interface
became \emph{internal}, which is what it will always be in a real system, and
the only signals reaching pins are switches, buttons and LEDs --- each one
flip-flop away from its pad.

That is not a workaround. It is the correct topology for the block, and it is
why the pad-to-pad and clock-insertion problems disappear rather than being
suppressed.

Final state: \textbf{WNS $+1.610$\,ns, hold $+0.122$\,ns, 0 of 277 failing
endpoints}, implied $F_{\max}\approx119$\,MHz, zero DRC violations, 320 of 320
nets routed, bitstream generated cleanly.

% =============================================================================
\section{Practices That Made the Difference}
% =============================================================================

\begin{enumerate}[leftmargin=1.6em,itemsep=5pt]

\item \textbf{Read individual paths, not the summary.} ``WNS $=-12.4$\,ns''
      says nothing about cause. The \sig{Source}/\sig{Destination} fields
      classify a path as port-to-port, port-to-flop, flop-to-port or
      flop-to-flop --- and that classification immediately showed the fault lay
      outside the logic.

\item \textbf{Read the implementation log.} The pblock failure produced no
      error and no warning. It appeared only as an \sig{INFO} line explaining
      that Vivado had rewritten the constraint. A clean run is not a correct
      run.

\item \textbf{Use the logic/route split as the stopping criterion.} At 87\,\%
      routing, logic optimisation is nearly futile. Knowing when to stop
      optimising is worth as much as knowing how to start.

\item \textbf{Make falsifiable predictions before each run.} Each change was
      preceded by an explicit prediction of the resulting WNS and failing
      endpoint count. When a prediction held, confidence in the model was
      earned rather than assumed.

\item \textbf{Take failed predictions seriously.} One prediction --- that a
      path containing no LUTs could not change --- was wrong: its \emph{routing}
      changed because a smaller netlist placed differently. The logic delay was
      indeed identical to three decimals, but the conclusion drawn from it was
      too strong. Recording this is more useful than recording the successes.

\item \textbf{Change one variable at a time.} Constraints, then abstraction,
      then RTL. Comparing an OOC result before a change against a pad-level
      result after it would have attributed the delta to the wrong cause.

\item \textbf{Test the noise hypothesis before believing it.} A $0.225$\,ns gain
      was initially dismissed as placement variance. Re-running the unchanged
      design produced a byte-identical result, proving the tool deterministic
      and the gain real.

\item \textbf{Choose the abstraction that matches the object.} Out-of-context
      synthesis is not a way to make numbers look better; it is the correct
      model for a block that will be integrated, and the pad-level run was the
      inaccurate one.

\item \textbf{Simulate before touching hardware.} Debugging a tick generator, a
      priority latch, a debouncer and a four-instance FSM through four LEDs is
      far slower than debugging them in a waveform.

\item \textbf{Prefer properties to vectors.} ``Every requesting master is served
      within $2N$ slots'' is checkable continuously across every scenario and
      survives changes in pointer phase; a hand-written expected sequence does
      neither.

\item \textbf{Synchronise every asynchronous input, but debounce only edges.}
      Metastability affects all asynchronous inputs; bounce only matters where a
      transition carries meaning. Conflating the two produces either missing
      synchronisers or unnecessary logic.

\item \textbf{Honour the module's stated contract at its boundary.} The arbiter
      requires \sig{REQ}/\sig{PRI} stable within a transfer, so the wrapper
      latches the priority selector at slot boundaries and the master model
      holds \sig{REQ} until completion. A raw switch would have violated the
      contract silently.

\item \textbf{State limitations precisely.} ``Starvation-free in slot count;
      latency bounded only under bounded transfer length; an indefinitely held
      lock is unbounded'' is a stronger position than an unqualified claim of
      fairness.

\end{enumerate}

% =============================================================================
\section{Summary}
% =============================================================================

\begin{center}
\begin{tabular}{@{}ll@{}}
\toprule
\textbf{Item} & \textbf{Status} \\
\midrule
Arbiter core, out-of-context timing & $+0.421$\,ns WNS, 67 LUTs, 19 FF \\
RTL simulation (wrapper)            & 88 slots, 9 scenarios, 0 errors \\
Post-implementation simulation      & 88 slots, netlist matches RTL exactly \\
Implementation (board wrapper)      & $+1.610$\,ns WNS, 0 / 277 failing \\
Hold                                & $+0.122$\,ns, 0 / 277 failing \\
DRC / routing                       & 0 violations; 320 / 320 nets routed \\
Bitstream                           & generated, 0 warnings \\
Hardware validation                 & \textbf{validated on Basys 3} (\sig{xc7a35t}) \\
\bottomrule
\end{tabular}
\end{center}

\bigskip\noindent
The defensible claim is therefore: \emph{an AMBA-style four-master arbiter with
round-robin and priority slot alternation, nested tie-breaking, lock override
and abandoned-grant recovery; verified in RTL and post-implementation
simulation; timing-closed at 100\,MHz with 1.61\,ns of margin; and validated on
hardware on a Basys 3 (\sig{xc7a35tcpg236-1}), where every documented behaviour
--- round-robin rotation, priority preemption, tie-break rotation within the
tied subset, lock hold with frozen alternation, and grace-window recovery from
an abandoned grant --- was observed on the board.}

\medskip\noindent
Board testing also produced a result worth recording, because it initially
looked like a fault. With all priorities equal and all four masters requesting,
the observed grant order was \texttt{M0 M3 M1 M0 M2 M1 M3 M2 \ldots}, which is
plainly not a single round-robin cycle. Separating it by \sig{slot_type} resolves
it: the round-robin slots ran \texttt{M0 M1 M2 M3 M0 M1 M2} and the priority
slots ran \texttt{M3 M0 M1 M2 M3 M0} --- two strict cycles, interleaved, offset
by three. The offset is itself evidence of correctness: \sig{tie_ptr} advances
only when a tie was actually resolved, so running an earlier preset with a
unique maximum left it stationary while \sig{rr_ptr} continued. That is exactly
the state-independence the specification requires, and it is visible from the
outside only if one remembers there are two pointers rather than one.

\medskip\noindent
Every clause is backed by a report. The more instructive half of the work,
however, was the diagnosis: roughly 95\,\% of an apparently catastrophic
$-12.4$\,ns violation belonged to the I/O ring and to a placement constraint
that had silently never applied --- and the arbiter's own logic, all eight
levels of it, contributed under one nanosecond throughout.

\end{document}
